\documentclass[10pt,a4paper]{article}

\usepackage{../ise-doc}
\usepackage{booktabs}

\title{BUGLEX --- Manual}
\author{Siddharth Shringarpure}
\date{}

\begin{document}
\singlespacing
\pagestyle{plain}
\maketitle

\section{Purpose}

This document explains how to use the bug report classification tool. The tool evaluates several text classification models for binary issue classification over five datasets: \texttt{caffe}, \texttt{incubator-mxnet}, \texttt{keras},
\texttt{pytorch}, and \texttt{tensorflow}. The main reported model is a feature-level hybrid that combines TF-IDF features with sentence embeddings.


\section{Getting Started}

Before using the tool, ensure your environment meets the system and software requirements detailed in \texttt{requirements.pdf}. This includes installing the \texttt{uv} package manager and ensuring a compatible \LaTeX{} toolchain is available if you intend to rebuild the report.

From the repository root, initialise the virtual environment with:
\begin{verbatim}
      uv sync
\end{verbatim}

Refer to \texttt{requirements.pdf} for platform-specific installation
instructions and a list of all Python dependencies.


\section{Main Commands}

\subsection{Run the Full Default Experiment}
To run the main experiment pipeline across all datasets and regenerate the main
report plots:

\begin{verbatim}
      uv run python3.13 src/run_experiments.py --with-plots
\end{verbatim}
This runs the default preprocessing mode, the main model comparison, the
embedding-dimension ablation, the plotting step, and all statistical tests (Wilcoxon and Friedman) automatically.


\subsection{Run Every Preprocessing Ablation}
To run all preprocessing variants as well as the default pipeline:

\begin{verbatim}
      uv run python3.13 src/run_experiments.py --all-preprocessing --with-plots
\end{verbatim}
This repeats the experiment workflow for:

\begin{itemize}
\item \texttt{none}
\item \texttt{stopwords\_all}
\item \texttt{stopwords\_keep\_negation}
\item \texttt{lemmatize}
\item \texttt{stopwords\_keep\_negation+lemmatize}
\end{itemize}


\subsection{Run One Dataset Only}
For a quicker or targeted run on a single dataset:

\begin{verbatim}
      uv run python3.13 src/run_experiments.py --dataset caffe
\end{verbatim}


\subsection{Run the Baseline Only}
To execute only the Naïve Bayes + TF-IDF baseline:

\begin{verbatim}
      uv run python3.13 src/run_baseline.py --dataset caffe
\end{verbatim}


\subsection{Regenerate Plots}
If the result CSV files already exist and only the figures need to be rebuilt:

\begin{verbatim}
      uv run python3.13 src/plot_results.py
\end{verbatim}


\subsection{Rebuild the Report}
To regenerate the \LaTeX{} table fragments and rebuild the main report PDF:

\begin{verbatim}
      uv run python3.13 src/tools/build_docs.py --report-only
\end{verbatim}
This command regenerates the \LaTeX{} table fragments from the current
\texttt{results/} CSV files and recompiles \texttt{docs/report/report.tex}.


\section{Outputs}

The main outputs are written to \texttt{results/}. Important files include:
\begin{itemize}
\item per-dataset run outputs such as
      \texttt{results/caffe\_all\_models\_768.csv}
\item top-level summary files such as
      \texttt{results/summary\_768.csv}
\item statistical summary files such as
      \texttt{results/wilcoxon\_summary\_768.csv}
\item ablation summaries such as
      \texttt{results/embedding\_ablation\_summary\_512\_256\_128\_64.csv}
\item report figures in \texttt{results/figures/}
\item generated report notes in \texttt{results/report\_notes.txt}
\end{itemize}

Cached embeddings are written to \texttt{results/embeddings/}. These caches are separated by dataset, model, dimension, and preprocessing mode, so routine reruns should not require manual cache deletion.


\section{Typical Usage}

For normal use, the simplest sequence is:

\begin{enumerate}
      \item ensure dependencies are installed via \texttt{uv sync}
      \item run the default experiment with
            \texttt{uv run python3.13 src/run\_experiments.py -{}-with-plots}
      \item if needed, run the full preprocessing ablation with
            \texttt{uv run python3.13 src/run\_experiments.py --all-preprocessing -{}-with-plots}
      \item rebuild the report with
            \texttt{uv run python3.13 src/tools/build\_docs.py -{}-report-only}
\end{enumerate}


\section{Troubleshooting}

\begin{itemize}
      \item If the first embedding run is slow, this is expected. The model is being downloaded and cached.
      \item If \texttt{-{}-help} feels slow, use the updated runner scripts. Their heavier imports are deferred until after argument parsing.
      \item If environment setup or report building fails, refer to \texttt{requirements.pdf} to ensure all dependencies and system tools are correctly installed.
\end{itemize}


\section{Artefact Location}

\texttt{manual.pdf} is placed in the repository root. The source for this PDF is stored in \texttt{docs/manual/}.

\end{document}
